Twenty citations were stated from standing knowledge --- author, year, venue --- pending
verification against the published record. Checking all twenty found \textbf{five errors}:

\begin{longtable}[]{@{}ll@{}}
\toprule
\begin{minipage}[b]{0.47\columnwidth}\raggedright
citation\strut
\end{minipage} & \begin{minipage}[b]{0.47\columnwidth}\raggedright
the error\strut
\end{minipage}\tabularnewline
\midrule
\endhead
\begin{minipage}[t]{0.47\columnwidth}\raggedright
Fisher (1935)\strut
\end{minipage} & \begin{minipage}[t]{0.47\columnwidth}\raggedright
given as ``ch.~21''; the randomization test is \textbf{\S21 within Chapter III}\strut
\end{minipage}\tabularnewline
\begin{minipage}[t]{0.47\columnwidth}\raggedright
Gelman \& Loken\strut
\end{minipage} & \begin{minipage}[t]{0.47\columnwidth}\raggedright
cited as a 2013 working paper; published as \textbf{2014, \emph{American Scientist} 102(6)}\strut
\end{minipage}\tabularnewline
\begin{minipage}[t]{0.47\columnwidth}\raggedright
Madani et al.~(2004)\strut
\end{minipage} & \begin{minipage}[t]{0.47\columnwidth}\raggedright
wrong authors --- and \textbf{the claim about it was backwards}\strut
\end{minipage}\tabularnewline
\begin{minipage}[t]{0.47\columnwidth}\raggedright
Taylor (2019)\strut
\end{minipage} & \begin{minipage}[t]{0.47\columnwidth}\raggedright
described as a quantile-loss neural network; it is \textbf{semiparametric asymmetric-Laplace}\strut
\end{minipage}\tabularnewline
\begin{minipage}[t]{0.47\columnwidth}\raggedright
Lin et al.\strut
\end{minipage} & \begin{minipage}[t]{0.47\columnwidth}\raggedright
dated 2023; published \textbf{2024}, \emph{Scientometrics} 129(4)\strut
\end{minipage}\tabularnewline
\bottomrule
\end{longtable}

\textbf{Four of the five were inherited from other people's descriptions of a source, or from a
metadata string, rather than from the source itself.} A citation copied from a citation is not
a verified citation, and five in twenty is higher than we would have guessed before measuring it.

The Madani case is the one worth dwelling on. It had been recorded, from a mention inside another
paper, as a probable early instance of stability being used as a quality signal --- that is, as
\emph{evidence for this paper's argument}. Reading it showed the opposite: they derive bounds relating
prediction disagreement to generalization error rather than assuming the relationship. \textbf{Left
unchecked it would have entered the paper as a fabricated ally.}

A coda, reported without irony. The Lin et al.~article --- the one whose date we got wrong --- is a
bibliometric study whose authors follow a method of their own devising called \emph{precise citation},
having previously published a paper on citation accuracy. \textbf{It carries a published correction
stating that its own citation information for Table 1 was incorrect.} If a paper about citation
accuracy, by authors who wrote a paper about citation accuracy, needs a citation correction, then
five in twenty is not a story about anyone being careless. It is a base rate, and the only
defence against a base rate is a procedure rather than an intention --- the same conclusion \Cref{sec:tautology}
reaches by a different route.

We resist turning this into a fifth finding. Two of the five errors ran in the direction of our
argument and three were neutral; n = 2 supports nothing, by the standard applied in \Cref{sec:power}. The
full bibliography, with every entry marked verified and every correction recorded, is
\texttt{paper/references.md}: \textbf{47 verified, 0 stated from memory, 0 placeholders.}
